\documentclass[11pt]{article}
\usepackage[margin=1in]{geometry}
\usepackage{amsmath,amssymb,amsthm,mathtools}
\usepackage[T1]{fontenc}
\usepackage[expansion=false]{microtype}
\usepackage[hidelinks]{hyperref}
\newtheorem{theorem}{Theorem}[section]
\newtheorem{lemma}[theorem]{Lemma}
\newtheorem{proposition}[theorem]{Proposition}
\newtheorem{remark}[theorem]{Remark}
\newcommand{\ddc}{\mathrm{d}\mathrm{d}^c}
\newcommand{\Ga}{\mathcal G_h}
\newcommand{\eps}{\varepsilon}
\newcommand{\norm}[1]{\left\|#1\right\|}
\title{Local critical bounds for exact Hermitian adjoints}
\author{}
\date{}
\begin{document}
\maketitle
\begin{abstract}
On a fixed compact Hermitian manifold, an explicit small ball of smooth
logarithmic densities in the $C^{2,\alpha}$ topology has uniformly bounded
densities for its exact Gauduchon maximizers. We estimate both the determinant
density and the associated adjoint probability density in $L^\infty$, hence
at the critical exponent $n/(n-1)$. The construction applies when the global
lower Monge--Amp\`ere mass is zero. A nonlinear contraction produces solution
forms close to the background, and a second contraction controls their
positive Gauduchon factors. Matrix arithmetic--geometric mean proves exact
maximality against every current test and identifies all equality cases.
The probability-normalized density family is infinite dimensional and
contains a path of distinct data with distinct exact maximizers. The result
is local in a strong density topology; weaker or larger bounded families
remain outside its scope.
\end{abstract}

\section{The two densities and the result}

Let $X$ be a connected compact complex manifold of dimension $n\ge3$,
let $h$ be a smooth positive Hermitian form, and let $\mathrm dV$ be a
smooth positive probability volume. We use
\[
 \ddc=\frac{i}{2\pi}\partial\bar\partial
\]
and exterior powers without divided factorials. Fix $0<\alpha<1$ and
H\"older norms from a finite smooth atlas and background connections.
Put
\[
 H=\frac{h^n}{\mathrm dV},\qquad V=\int_Xh^n,
 \qquad F_* =\log(H/V).
\]
Thus $e^{F_*}\mathrm dV=h^n/V$ is a probability measure.

The compact test class consists of positive $(n-1,n-1)$-currents
\[
 \Ga=\left\{T:\ddc T=0,\ \int_Xh\wedge T=1\right\}.
\]
In an $h$-unitary frame set $\sigma_T=h\wedge T$ and use the positive
Hermitian matrix density $P_T$ characterized by
\[
 \beta\wedge T=\frac1n\operatorname{tr}(BP_T)\,\sigma_T,
 \qquad \operatorname{tr}P_T=n,
\]
where $B$ is the matrix of $\beta$. For real $F$ define
\[
 \Gamma_F(T)=\left(\frac{e^F\mathrm dV}{h^n}\right)^{1/n}
                  \det(P_T)^{1/n}\sigma_T,
 \qquad M(F)=\max_{T\in\Ga}\Gamma_F(T)(X).
\]
At an exact maximizer, its determinant-weighted probability is
$\nu_{T,F}=\Gamma_F(T)/M(F)$. We write $q_T=\mathrm d\nu_{T,F}/\mathrm dV$.
For a smooth positive $T=\gamma^{n-1}$ its determinant density is
\[
 k_T=\left(\frac{\gamma^n}{\mathrm dV}\right)^{(n-1)/n}.
\]
The identity $M(F)=e^{-B_0(F)/n}$ identifies $\nu_{T,F}$ with an adjoint
probability of the functional $B_0$ that normalizes the equation. In this paper
exact maximality and both density formulas will be proved directly for
the constructed family.

\begin{theorem}\label{thm:endpoint}
There is an explicit $\eps>0$, defined in Section~\ref{sec:constants}
using only the fixed background and linear elliptic inverse norms, with
the following property. For
\[
 \mathcal F_\eps=\left\{F\in C^\infty(X,\mathbb R):
       \int_Xe^F\mathrm dV=1,\quad
       \norm{F-F_*}_{C^{2,\alpha}}\le\eps\right\},
\]
every exact maximizing current is the same smooth strictly positive
current $T_F$. Its two densities satisfy
\[
 \norm{q_{T_F}}_{L^\infty(\mathrm dV)}\le C_\nu,
 \qquad \norm{k_{T_F}}_{L^\infty(\mathrm dV)}\le C_k,
\]
where, for the normalized baseline factor $\chi_0$ constructed below,
\[
 a=\max_X\chi_0,\qquad H_{\max}=\max_XH,
 \qquad C_\nu=\frac{3^{n+1}}4aH_{\max},
 \qquad C_k=\frac{3^n}2aH_{\max}^{(n-1)/n}.
\]
Consequently, at $p=n/(n-1)$,
\[
 \norm{q_{T_F}}_{L^p(\mathrm dV)}\le C_\nu^{1/n},
 \qquad \norm{k_{T_F}}_{L^p(\mathrm dV)}\le C_k.
\]
The family is infinite dimensional and contains a path of distinct data
and distinct exact maximizing currents. The theorem holds on the same
fixed $h$ when $v_-(h)=0$; no positive lower mass is assumed.
\end{theorem}

All data here are smooth. Their boundedness is in the specified
$C^{2,\alpha}$ norm, with no uniform bounds on higher derivatives.
The estimates also hold on the entire small smooth ball if the
probability-normalization condition is omitted.

\section{Linear inverses and a baseline factor}\label{sec:linear}

We use real scalar elliptic theory on a compact manifold without boundary.
For clarity, the passage to the full H\"older spaces is as follows.
For a smooth elliptic second-order operator $P$, a two-sided parametrix
$E$ has order $-2$ and
\[
 EP=I-S,\qquad PE=I-R,
\]
where $R,S$ have smooth kernels. The construction and distributional
regularity are given in \cite[Chapter 7, Sections 4 and 10]{Taylor7}.
The H\"older mapping property is
$E:C^{s,\alpha}\longrightarrow C^{s+2,\alpha}$;
see \cite[Chapter 13, Section 8, particularly (8.7), (8.17) and
(8.22)]{Taylor13}. Consequently $v=EPv+Sv$ gives
\[
 \norm{v}_{C^{s+2,\alpha}}
  \le C\bigl(\norm{Pv}_{C^{s,\alpha}}+\norm{v}_{C^0}\bigr),
\]
for each nonnegative integer $s$. A distribution with $Pv$ smooth is
smooth by the same parametrix identity.

The Fredholm alternative in \cite[Chapter 7, (10.4)--(10.6)]{Taylor7}
first applies on $H^2\to L^2$. If $f\in C^{s,\alpha}$ is orthogonal
to all smooth formal-adjoint kernel elements, it therefore has an
$H^2$ solution. The parametrix identity upgrades that solution to
$C^{s+2,\alpha}$. Conversely, integration by parts gives the same
orthogonality condition. This identifies the H\"older range as the
closed intersection of finitely many smooth integral constraints.
Its kernel is smooth and finite dimensional, so the H\"older operator
is Fredholm, with the same kernel, cokernel and index as the Sobolev
operator. This argument does not require smooth functions to be dense
in the full H\"older norm. Index invariance follows from
\cite[Section 7, Proposition 7.4 and its proof]{TaylorFA}.
A real Laplacian has index zero by its self-adjointness and the same
smooth kernel--cokernel description. Changing the reference volume
changes the adjoint representatives by a smooth positive factor,
without changing the dimension or compatibility condition.

We also recall the strong maximum principle in the form needed here.
If $P=a^{ij}\partial_i\partial_j+b^i\partial_i+c$ is locally uniformly
elliptic, $c\le0$, $Pv\ge0$ and $v$ attains a nonnegative interior
maximum $M$, put $w=M-v\ge0$. Then $Pw=cM-Pv\le0$.
If $w$ were positive somewhere and zero somewhere in the connected
domain, a coordinate ball $B_R(y)$ in its positive set could touch
its zero set at an interior point $z$. On the annulus
$R/2<|x-y|<R$, the barrier
\[
 b_\kappa(x)=e^{-\kappa|x-y|^2}-e^{-\kappa R^2}
\]
satisfies $Pb_\kappa>0$ for sufficiently large $\kappa$: its positive
$4\kappa^2a^{ij}(x_i-y_i)(x_j-y_j)$ term dominates the bounded
terms of order $\kappa$ and order zero. Choose $\tau>0$ so that
$w-\tau b_\kappa\ge0$ on both boundary spheres. A negative interior
minimum would give $P(w-\tau b_\kappa)\ge0$, contradicting its strict
negativity. Thus $w\ge\tau b_\kappa$ on the annulus. Its inward
derivative at $z$ is then positive, contradicting $\nabla w(z)=0$.
This proves that $v$ is constant. Applying the result to $-v$ gives
the strong minimum principle used below. The argument is local and
then propagates through connected coordinate neighborhoods.

Define
\[
 L_hv=\frac{n\ddc v\wedge h^{n-1}}{h^n}.
\]
Its coefficients are smooth, its principal part is uniformly elliptic,
and it has no zeroth-order term. The maximum principle and connectedness
give $\ker L_h=\mathbb R$. Deform its principal coefficients through
positive definite coefficients to a real Laplacian and its lower-order
coefficients to those of that Laplacian. The index is therefore zero,
so its formal adjoint relative to $\mathrm dV$ has a one-dimensional
smooth kernel. Let $p_0$ be a nonzero real generator.

The generator has one sign. Indeed, the range of $L_h$ contains no
nonzero nonnegative smooth function: $L_hv\ge0$ on the compact manifold
forces $v$ constant. If $p_0$ took both signs, nonnegative smooth bumps
supported in its positive and negative sets would give a nonzero
nonnegative $g$ with $\int_Xgp_0\mathrm dV=0$. Fredholm compatibility
would put $g$ in the range, a contradiction. Choose $p_0\ge0$.
It is strictly positive. In coordinates the adjoint has a zeroth-order
coefficient $c_0$; subtract a constant $C\ge\max\{0,\sup c_0\}$.
The resulting operator $P$ has nonpositive zeroth-order coefficient
and $Pp_0=-Cp_0\le0$. An interior zero would force $p_0\equiv0$
by the strong minimum principle. Normalize $\int_Xp_0\mathrm dV=1$.

Set $\chi_0=p_0/H$. Integration by parts gives
\[
 \int_XL_hv\,p_0\mathrm dV
       =n\int_Xv\,\ddc(\chi_0h^{n-1}),
\]
and hence
\begin{equation}\label{eq:baseline}
 \ddc(\chi_0h^{n-1})=0,\qquad
 \int_X\chi_0h^n=1,\qquad \chi_0>0.
\end{equation}
Write $\mathrm d\mu_0=\chi_0h^n$.

Let $E_4=\{v\in C^{4,\alpha}:\int_Xv\mathrm dV=0\}$. The map
\[
 \mathcal L:E_4\times\mathbb R\longrightarrow C^{2,\alpha},
 \qquad \mathcal L(v,c)=L_hv-c
\]
is an isomorphism. For prescribed $g$ its scalar is
$c=-\int_Xg\mathrm d\mu_0$, and the remaining equation is
$L_hv=g-\int_Xg\mathrm d\mu_0$. Compatibility and the normalization
give existence and uniqueness. Fix $K\ge1$ bounding its inverse for
the product norm $\norm{v}_{C^{4,\alpha}}+|c|$.

For a positive smooth $(1,1)$-form $\omega$, define
\[
 A_\omega\eta=\frac{\ddc(\eta\omega^{n-1})}{\mathrm dV}.
\]
The formal adjoint of $A_h$ is
$\eta\mapsto(\ddc\eta\wedge h^{n-1})/\mathrm dV=(H/n)L_h\eta$.
It has constant kernel and index zero. Thus $A_h$ has kernel spanned
by $\chi_0$ and range
$Z_0=\{g\in C^{0,\alpha}:\int_Xg\mathrm dV=0\}$.
Its restriction to
\[
 E_\chi=\left\{\eta\in C^{2,\alpha}:\int_X\eta h^n=0\right\}
\]
is an isomorphism onto $Z_0$. Denote its inverse by $\mathcal R$ and
fix $G\ge1$ bounding $\norm{\mathcal R}$.

\section{Explicit radius and two contractions}\label{sec:constants}

Rescale the fixed form and endomorphism H\"older norms so that
endomorphism composition is submultiplicative and the norms dominate
pointwise operator norms. Choose $D\ge1$ with
\[
 \norm{\ddc v}_{C^{2,\alpha}}\le D\norm{v}_{C^{4,\alpha}},
 \qquad
 \norm{h^{-1}\ddc v}_{C^{2,\alpha},\mathrm{alg}}
                          \le D\norm{v}_{C^{4,\alpha}}.
\]
Let $C_{\mathrm{tr}}\ge1$ bound the trace map and put
$Q=2C_{\mathrm{tr}}D^2$. For $\norm{v}_{C^{4,\alpha}}\le1/(2D)$ set
\[
 N(v)=\log\det(I+h^{-1}\ddc v)-L_hv.
\]
The Banach algebra series
\[
 N(v)=\sum_{j\ge2}\frac{(-1)^{j+1}}j
                  \operatorname{tr}\bigl((h^{-1}\ddc v)^j\bigr)
\]
and the telescoping formula for a difference of matrix powers imply
\begin{align}
 \norm{N(v)}_{C^{2,\alpha}}&\le Q\norm{v}_{C^{4,\alpha}}^2,\label{eq:remainder}\\
 \norm{N(v)-N(w)}_{C^{2,\alpha}}
  &\le Q\bigl(\norm{v}_{C^{4,\alpha}}+
          \norm{w}_{C^{4,\alpha}}\bigr)\norm{v-w}_{C^{4,\alpha}}.
 \nonumber
\end{align}
Indeed, if $\rho\le1/2$ bounds both matrix norms, the difference
series is bounded by
$C_{\mathrm{tr}}\rho(1-\rho)^{-1}$ times the matrix difference.

Choose $C_A\ge1$ so that, for $\norm{\omega-h}_{C^{2,\alpha}}\le1$,
\[
 \norm{A_\omega-A_h}_{C^{2,\alpha}\to C^{0,\alpha}}
                         \le C_A\norm{\omega-h}_{C^{2,\alpha}}.
\]
This is a finite polynomial coefficient estimate: expand
$\omega^{n-1}-h^{n-1}$, multiply by $\eta$ and apply $\ddc$.
It uses at most two derivatives of $\omega$.

All constants so far are fixed before solving the equation. Define
\begin{equation}\label{eq:radius}
 m=\min_X\chi_0>0,\qquad C_0=\norm{\chi_0}_{C^{2,\alpha}},\qquad
 r=\min\left\{\frac1{2D},\frac1{4KQ},
                   \frac{m}{4GC_ADC_0}\right\},\qquad
 \eps=\frac r{2K}.
\end{equation}
In particular $m\le C_0$. There is no lower mass or inverse lower
mass in these formulas.

For $f=F-F_*$ consider the map on the closed radius-$r$ ball in
$E_4\times\mathbb R$
\[
 (v,c)\longmapsto\mathcal L^{-1}\bigl(f-N(v)\bigr).
\]
When $\norm f_{C^{2,\alpha}}\le\eps$, its image has norm at most
$K\eps+KQr^2\le3r/4$ and its Lipschitz constant is at most
$2KQr\le1/2$. Contraction therefore gives $(u,c)$ with
\[
 \norm u_{C^{4,\alpha}}+|c|\le r,\qquad
 \log(\omega^n/h^n)-c=f,\qquad \omega=h+\ddc u.
\]
The pointwise norm bound $\norm{h^{-1}\ddc u}\le Dr\le1/2$ gives
\begin{equation}\label{eq:metric-comparison}
 \tfrac12h\le\omega\le\tfrac32h,\qquad
 \omega^n=e^{F+b}\mathrm dV,\quad b=\log V+c.
\end{equation}
For smooth $F$, $u$ is smooth. Differentiate the local log-determinant
equation; its linearization has uniformly positive principal
coefficients $\omega^{-1}$, initially in $C^{2,\alpha}$.
Interior Schauder estimates on the differentiated equations improve
the regularity one order at a time. Equivalently, the local nonlinear
elliptic bootstrap in \cite[Proposition and complete proof]{TaylorReg}
applies with order two, initial regularity $C^{4,\alpha}$, smooth
right side and positive definite linearization. It gives all
derivatives. Uniform higher derivative bounds on $F$ are not used.

Now set $\Delta A=A_\omega-A_h$. Each $A_\omega$ has integral zero,
so the following right side belongs to $Z_0$. On $E_\chi$ solve
\[
 \eta=-\mathcal R\Delta A(\chi_0+\eta).
\]
It is an affine contraction on this complete space, since
$G\norm{\Delta A}\le GC_ADr\le m/(4C_0)\le1/4$.
Its unique fixed point satisfies
\[
 \norm\eta_{C^{2,\alpha}}
  \le\frac{GC_ADrC_0}{1-GC_ADr}\le2GC_ADrC_0\le m/2.
\]
Thus $\chi=\chi_0+\eta$ obeys
\begin{equation}\label{eq:factor}
 \ddc(\chi\omega^{n-1})=0,\quad \int_X\chi h^n=1,
 \quad \chi\ge m/2,
 \quad \chi\le\chi_0+m/2\le\tfrac32\max_X\chi_0.
\end{equation}
It is smooth by elliptic regularity. Every $C^{2,\alpha}$ solution
with this normalization satisfies the same affine equation, so it
equals $\chi$.

\section{Exact maximizers and both uniform bounds}

Define
\[
 Z=\int_X\chi h\wedge\omega^{n-1},\qquad
 T_F=\frac\chi Z\omega^{n-1}.
\]
It is a smooth strictly positive member of $\Ga$. Integration by parts
and \eqref{eq:factor} give
\[
 \int_X\chi(\omega-h)\wedge\omega^{n-1}
       =\int_Xu\ddc(\chi\omega^{n-1})=0,
 \qquad Z=\int_X\chi\omega^n.
\]
Diagonalizing $\omega$ relative to $h$, the ratio
$h\wedge\omega^{n-1}/h^n$ is the average of products of $n-1$
eigenvalues. Consequently
\begin{equation}\label{eq:Z}
 2^{-(n-1)}\le Z\le(3/2)^{n-1}.
\end{equation}

For any $S\in\Ga$, smoothness of $u$ and $\ddc S=0$ imply
$\int_X\omega\wedge S=1$. Let $W$ be the matrix of $\omega$ in the
$h$-unitary frame and put $P=P_S$. By \eqref{eq:metric-comparison},
$e^F\mathrm dV/h^n=e^{-b}\det W$. Matrix arithmetic--geometric mean gives
\[
 \det(W)^{1/n}\det(P)^{1/n}
 =\det(W^{1/2}PW^{1/2})^{1/n}\le\frac1n\operatorname{tr}(WP).
\]
It follows as an inequality of measures that
\[
 \Gamma_F(S)\le e^{-b/n}\omega\wedge S,
 \qquad \Gamma_F(S)(X)\le e^{-b/n}.
\]
This includes singular test currents. For $T=t\omega^{n-1}$ with a
smooth positive factor $t$, its trace-normalized matrix is
$P_T=nW^{-1}/\operatorname{tr}(W^{-1})$. Thus equality holds for $T_F$:
\begin{equation}\label{eq:exact}
 M(F)=e^{-b/n},\qquad
 \Gamma_F(T_F)=e^{-b/n}\omega\wedge T_F.
\end{equation}
In particular $b=B_0(F)$ by the stated normalizing-functional identity.

Every exact maximizer equals $T_F$. Equality of the total integrals
forces the nonnegative deficit measure to be zero. Equality in matrix
arithmetic--geometric mean then gives, $\sigma_S$-almost everywhere,
\[
 W^{1/2}P_SW^{1/2}=\lambda I,\qquad
 \lambda=\frac n{\operatorname{tr}(W^{-1})}>0.
\]
Hence $S=\rho\omega^{n-1}$ for a nonnegative scalar measure coefficient
$\rho$. The distributional equation $\ddc S=0$ is $A_\omega\rho=0$.
Smooth-coefficient elliptic regularity makes $\rho$ smooth. Normalize
it by its positive finite integral against $h^n$. The uniqueness in
\eqref{eq:factor} identifies it with $\chi$, and the $h$-trace
normalization fixes its remaining scale to $1/Z$. This proves
uniqueness by the equality equation and ellipticity.

The probability and determinant densities are now explicit:
\begin{equation}\label{eq:densities}
 \nu_{T_F,F}=\omega\wedge T_F=\frac\chi Z\omega^n,
 \qquad q_F=\frac\chi Z\frac{\omega^n}{\mathrm dV},\qquad
 k_{T_F}=\frac\chi Z
                 \left(\frac{\omega^n}{\mathrm dV}\right)^{(n-1)/n}.
\end{equation}
For the last identity take
$\gamma=(\chi/Z)^{1/(n-1)}\omega$, so $T_F=\gamma^{n-1}$.
Equivalently $k_{T_F}=e^{-F/n}M(F)q_F$.
The adjoint and supporting properties can also be checked directly.
For every smooth $v$,
\[
 \int_X\frac{n\ddc v\wedge\omega^{n-1}}{\omega^n}
                  q_F\mathrm dV
       =\frac nZ\int_X\ddc v\wedge\chi\omega^{n-1}=0.
\]
For another continuous log-density $G$, testing its maximum at $T_F$
gives
\[
 M(G)\ge M(F)\int_Xe^{(G-F)/n}\mathrm d\nu_{T_F,F}.
\]
Jensen's inequality and $M=e^{-B_0/n}$ therefore give
$B_0(G)\le B_0(F)-\int_X(G-F)\mathrm d\nu_{T_F,F}$,
the required supporting inequality with its negative sign.
Equation \eqref{eq:Z} gives
\[
 \frac\chi Z\le\frac{3a}2\,2^{n-1},\qquad
 \frac{\omega^n}{\mathrm dV}\le(3/2)^nH_{\max}.
\]
Multiplication gives exactly $q_F\le C_\nu$ and $k_{T_F}\le C_k$
in Theorem~\ref{thm:endpoint}. Since $\mathrm dV$ is a probability volume,
$\norm{k_{T_F}}_{L^p}\le C_k$. Since $q_F$ itself integrates to one,
\[
 \int_Xq_F^p\mathrm dV\le C_\nu^{p-1},\qquad
 \norm{q_F}_{L^p}\le C_\nu^{(p-1)/p}=C_\nu^{1/n}.
\]

\section{Nontriviality at the fixed zero-mass background}

The radius in \eqref{eq:radius} is positive for every smooth positive
$h$, including the given one with $v_-(h)=0$. The background remains
fixed throughout the following density variations.

Choose a point $x_0$ and holomorphic coordinates unitary for $h$ there.
A cutoff equal to one near $x_0$, multiplied by
$|z_1|^2-|z_2|^2$, gives a global smooth real bump $\psi$ whose
complex Hessian matrix at $x_0$ is nonzero and non-scalar. Adjust its
$\mathrm dV$-mean to zero. For sufficiently small real $t$ put
\[
 \omega_t=h+t\ddc\psi,\qquad V_t=\int_X\omega_t^n,
 \qquad F_t=\log\frac{\omega_t^n}{V_t\mathrm dV}.
\]
These forms are positive, $F_t$ are smooth probability log-densities,
and $F_t\to F_*$ in $C^{2,\alpha}$. Shrink the interval so they lie
in $\mathcal F_\eps$ and the exact solution
$u=t\psi$, $b=\log V_t$ lies in the local contraction ball.

They are pairwise distinct. If $F_s=F_t$, then
$\log(\omega_s^n/\omega_t^n)=\log(V_s/V_t)$ is constant.
At a maximum and minimum of $(s-t)\psi$, the Hessian sign makes this
constant respectively nonpositive and nonnegative. Thus it is zero.
The mean-value formula for log determinants gives
\[
 \int_0^1\operatorname{tr}_{\omega_t+\tau(\omega_s-\omega_t)}
                          \ddc((s-t)\psi)\,\mathrm d\tau=0.
\]
This is a uniformly elliptic scalar equation without a zeroth-order
term. The maximum principle forces $(s-t)\psi$ constant, hence $s=t$.
Probability normalization also forbids nonzero constant shifts between
these data.

Their exact currents are distinct as well. If $T_s=T_t$, their positive
smooth factors imply that $\omega_s^{n-1}$ and $\omega_t^{n-1}$ are
proportional at $x_0$. The cofactor matrix identity implies that the
positive forms themselves are proportional. Their matrices there are
$I+sQ$ and $I+tQ$ with $Q$ non-scalar. Linear independence of $I,Q$
then forces $s=t$.

The entire family is infinite dimensional. Every sufficiently small
smooth perturbation $f$ can be normalized by
\[
 F=F_*+f-\log\int_Xe^{F_*+f}\mathrm dV.
\]
The derivative at zero is $f-\int_Xf e^{F_*}\mathrm dV$, whose only
constant-direction kernel is the constants. Arbitrarily many independent
nonconstant smooth directions remain inside $\mathcal F_\eps$.

The density hypothesis proved here is strong-topology smallness around
the explicit probability density $h^n/V$. It gives the requested
critical estimate at zero lower mass, and in fact an $L^\infty$ estimate.
It gives no estimate for arbitrary $C^0$-bounded or merely
$L^\infty$-bounded continuous log-densities, large $C^{2,\alpha}$
families, all normalized Gauduchon tests or near-maximizers. No general
singleton-support assertion is inferred from an endpoint estimate.
Those broader questions remain open.
\begin{thebibliography}{4}
\bibitem{Taylor7} M.~E.~Taylor,
\emph{Pseudodifferential Operators}, Chapter 7 of
\emph{Partial Differential Equations}, author-hosted chapter,
\url{https://mtaylor.web.unc.edu/wp-content/uploads/sites/16915/2018/04/chap7.pdf}.
\bibitem{Taylor13} M.~E.~Taylor,
\emph{Function Space and Operator Theory for Nonlinear Analysis},
Chapter 13 of \emph{Partial Differential Equations}, author-hosted chapter,
\url{https://mtaylor.web.unc.edu/wp-content/uploads/sites/16915/2018/04/chap13.pdf}.
\bibitem{TaylorFA} M.~E.~Taylor,
\emph{Functional Analysis}, Section 7,
author-hosted notes,
\url{https://mtaylor.web.unc.edu/wp-content/uploads/sites/16915/2018/04/fctnl.pdf}.
\bibitem{TaylorReg} M.~E.~Taylor,
\emph{Regularity of Solutions to Nonlinear Overdetermined Elliptic PDE},
author-hosted note,
\url{https://mtaylor.web.unc.edu/wp-content/uploads/sites/16915/2022/05/overellip.pdf}.
\end{thebibliography}
\end{document}
